\scanpage{95}
trouve l'isom
\[
\unsure{\kappa_\alpha|K_\alpha|k_\alpha}:T_0^{\otimes2}\simeq T_{\Sigma,t}
\]
\unsure{d'où|où} une courbe elliptique $E_\alpha$, \inserted{corr. à $\underline I,t$} munie d'un \deleted{\unsure{rigidification}} \inserted{épinglage} de Legendre subordonné à $\alpha$ (puisque $T_{E_\alpha,0}\simeq T_0$ \ldots).
% slanted
\margin{c'est la courbe elliptique \uline{universelle}, correspondant à $\underline I,t$ et munie d'un épinglage de Legendre subordonné à $\alpha$ \ldots}
Ceci posé, pour une courbe elliptique quelconque $E$ sur $S$, corr. à $\underline I,t$, les épinglages de Legendre sur $E$ subordonnés à $\alpha$ forment (on l'a vu) les sections d'un $\mu_2$-torseur $\mathscr L_\alpha(E)$ sur $S$, et on voit de suite qu'on a un isom de $\mu_2$-torseurs
\[
\deleted{\mathscr L_\alpha(E)\longleftarrow}
\]
\[
\begin{aligned}
\underbrace{\underline{\operatorname{Isom}}_{S,\underline I}(E_\alpha,E)}_{
\substack{\text{l'indice }\underline I\text{ signifie}\\
\text{qu'on prend les iso.}\\
\text{compatibles avec}\\
\underline I\simeq {}_2E\simeq {}_2E_\alpha}}
&\xrightarrow{\sim}\mathscr L_\alpha(E)\\
u&\longmapsto u(\kappa_\alpha)
\end{aligned}\tag*{(57)}
\]
Si, au lieu d'une section \inserted{$\alpha$} de $\mu_4^\times\mathbin{\wedge}\underline\omega_{\underline I}$, on disposait d'une section $\omega$ de $\underline\omega_{\underline I}$ lui-même, on pourrait encore définir une courbe elliptique $E_\omega$ canonique corr. à $\underline I,t$, et munie d'un ``ép. de Legendre ``naïf'''' \inserted{(\unsure{i.e.|lire} d'un iso. $T_{E,0}\simeq\underline O_\Sigma(1)(t)$ tel que \ldots)}, et
% archive sheet number: 94
